\section{Finite certificates and their verification}\label{app:verifier}

The verification programs, exact rational certificates and replay
instructions are in the
\href{https://raw.githubusercontent.com/royvanrijn/jacobian-research/b854dddcda88fd9c733d18e14a8ff5ae4a0b25b8/unpublished/eisenstein-restricted-prime-sums/source.zip}{frozen certificate supplement}
(GitHub commit \texttt{b854dddcda88}), under \texttt{certificates/}.
The archive also retains the sources of that edition; its certificate
files apply to this article. The separate TeX sources compile without
executing verification code. Table~\ref{tab:exact-coverage} gives the
complete finite domains and the interval and descent arguments joining
them to the infinite ranges.

\paragraph{Membership and coverage.}
The prime sieve deletes every multiple of a forbidden prime.
For the Eisenstein indexing \(6j+1\), any inert factor has another
inert factor counted with multiplicity, so one is below the square
root; primes below \(10000\) suffice in the stated domain.
The other three discriminants use the complete prime sieves specified
in their sections.

If \(B_C\) is the packed membership array, the unions of shifts
\(B_C\ll a\), with the indicated ramified dilation in a mixed family,
contain only genuine positive sums. Comparing the full truncated union
with the predicted target mask proves coverage of every eligible
integer and exclusion of every forbidden one. For a listed exception,
testing every possible summand up to half the target proves
nonrepresentation. The exact programs reconstruct these arrays and
target comparisons, not merely their hashes.

Independent coordinate enumeration authenticates the full membership
arrays at the certificate bounds. At \(-7,-11\), use
\(4Q=(2x+y)^2+|d|y^2\), sign reversal and \(x\mapsto-x-y\)
to restrict to \(y\ge0,\ 2x+y\ge0\), with integer square-root limits.
At \(-8\), independent sign changes restrict to the first quadrant.
Apply the coordinate gcd test throughout. At \(-11\), retain only
positive odd norms with \(11\nmid Q(x,y)\), thereby authenticating
\(C\), rather than \(P=C\mathbin{\dot\cup}11C\).
The full bytearrays are compared; deleting a unit, inserting a
forbidden prime power or corrupting a target mask must be rejected.

\paragraph{Exact arithmetic and completeness.}
Divisor maxima are exhausted by nonincreasing exponents on the
smallest eligible primes, with the explicit prime and exponent limits
proved in each section. Omitted Euler factors are at most one.
Rational rounding and arithmetic-progression tails bound the prime
sums. Integer comparisons after clearing positive denominators then
certify every target in each interval. The interval endpoints meet
the finite scans and infinite estimates; descent exhausts the
remaining targets.
In particular the three rows of \eqref{eq:m8-intervalcert} have
minimum margin greater than \(1/256\). Their exact minima are recorded
in \texttt{primitive\_norm\_minus8\_checks.json}, under
\texttt{interval\_certificate}, together with all \(107\) inequalities.

\paragraph{Receipt conventions.}
The supplement's four sumset verifiers and coordinate authenticator
produce the corresponding JSON receipts; the theta, Euler and local
checks are included there too. Membership arrays are ordered zero-one
bytes indexed by the integer, or by \(j\) for \(6j+1\).
Packed bit \(n\) is the coefficient of \(2^n\).
Eisenstein witness digests use the least short index at increasing
target indices, encoded as signed little-endian \(32\)-bit integers.
The ordinary and mixed arrays have their respective domain lengths.
The expected membership and witness digests are in
\texttt{binary\_certificates\_receipt.json}; the other families'
digests and rational margins are in their named JSON receipts.
The standalone verifiers use integer arithmetic and Python's standard
library; their assertions reject optimized execution.
Hashes identify the supplied bytes; the membership, interval and
descent arguments above establish their mathematical role.
